\FrontHeading{Abstract}
\noindent Retinopathy of prematurity (ROP) is a leading preventable cause of childhood blindness, and the limit on screening it is the number of clinicians able to read fundus photographs inside the short window in which treatment still works. This project built RetinAI, a software system that takes every photograph of one infant eye, grades each photograph into the six ICROP-3 classes (normal, stages~1--3, stage~4/5 and aggressive ROP), and returns one stage for the eye by the worst-view rule that an ophthalmologist applies when reading an examination. Our first version was a photograph-level EfficientNetV2-S classifier; the post-proposal review asked for a different architecture, eye-level input and output, higher accuracy and a publishable result. This report describes the second version and its evaluation.

Rather than tuning until a number improved, we fixed an evaluation protocol before any new model was trained: patient-grouped five-fold cross-validation over 6{,}700 photographs from 1{,}528 infants at three hospitals, a fourth hospital held out entirely (663 photographs, 230 eyes, 116 infants in the development half; a locked half of 661 photographs opened only once), one pre-registered primary endpoint (eye-level quadratic weighted kappa, QWK, on the locked half), infant-clustered bootstrap intervals, and a selection rule that never looks at the held-out hospital. Under this protocol we trained and compared nine models with two random seeds each: a retina-pretrained vision transformer (RETFound-Green, A), a natural-image-pretrained transformer of identical architecture (DINOv2 ViT-S/14, B), the EfficientNet baseline (C), B trained one examination session at a time on the Ostrava cohort (D), an ensemble (E), B adapted to the new hospital's unlabelled photographs (F), B at its native 518-pixel resolution (G), and, in a final pre-registered round (v4), B on a training pool cleaned by a per-camera frame-quality audit (H) and the same pool with a regularising training recipe (I). Every training run, every prediction and every score is committed to version control, and all 24 training sessions ran on free Kaggle GPUs, at about 33--35~minutes per fold at 392~pixels.

The measured results are: cross-validation QWK of 0.825 (95\% CI 0.737--0.911) for DINOv2, the model selected by the pre-registered rule, against 0.816 for the baseline; eye-level QWK at the unseen hospital of 0.756 (0.679--0.821) for DINOv2 and 0.760 for the baseline; accuracy 63.0\% per eye and 76.6\% per photograph; and any-ROP detection AUC of 0.937 per eye. The application serves DINOv2 as the mean of its five fold models, which scores \DeployedQWK{} (\DeployedQWKCI) per eye at the unseen hospital. No new model beats the baseline at the unseen hospital. The data audit and the recipe did not reach the pre-set adoption margin of $+0.010$ cross-validation QWK over B (H $-0.013$, I $-0.004$), so B stays selected; model~I scored best at the unseen hospital (0.776, paired gain over B $+0.020$, 95\% CI $+0.003$ to $+0.041$), a descriptive result below the 0.08 difference the protocol can detect. Two gains are real. Stage-1 recall at the new hospital rises from 0.40 (baseline) to 0.50--0.52 for transformers A, B, G, H and I, and training at the unit the labels describe (model~D) raises Ostrava examination-level QWK from 0.701 to 0.846 (+0.145, 95\% CI 0.094--0.257). Resolution, ensembling, ordinal decoding, unsupervised adaptation, data cleaning and the recipe do not help; the nine models lie within 0.032 QWK of one another, which places the ceiling in the labels, not the model.

The contributions are an eye-level ROP screening application with a PDF report, the first architecture-matched comparison of retinal versus natural-image self-supervised pre-training for multi-class ICROP-3 staging, a pre-registered data-audit and recipe round with a leave-one-hospital-out guard, a fully reproducible multi-hospital evaluation, and the finding that matching the training unit to the labelling unit is the one change that moves the result. The main limitations are that labels were not re-graded by a clinician, that the severe classes have too few eyes at the new hospital to evaluate, that the development half has now been examined for nine models, and that the locked half awaits its single approved scoring.

\noindent\textbf{Report layout}\par
\noindent Chapter~1 gives the clinical background, problem statement and objectives. Chapter~2 reviews the literature on automated ROP staging and identifies the gaps this work addresses. Chapter~3 states the requirements, constraints and risks. Chapter~4 describes the system architecture, the design alternatives, the model and the evaluation protocol. Chapter~5 covers the implementation: data pipeline, training on Kaggle, scoring code and the web application. Chapter~6 gives the testing methodology, datasets and metrics. Chapter~7 reports and discusses the measured results. Chapter~8 concludes against the objectives and sets out future work. Appendices give the protocol amendments, raw per-fold results, reproducibility information and project records.

\noindent\textbf{Keywords}\par
\noindent Retinopathy of prematurity; ICROP-3 staging; fundus photography; vision transformer; DINOv2; RETFound; transfer learning; data quality audit; external validation; pre-registration; quadratic weighted kappa; eye-level evaluation; medical image classification.
\clearpage
